\chapter{Reconstrucción fotorrealista mediante Gaussian Splatting}
\label{Reconstrucción fotorrealista}

Este anexo desarrolla el módulo de reconstrucción fotorrealista de la escena introducido en el
capítulo \ref{Solución}. Es una \textbf{contribución instrumental} del proyecto, no un objetivo del
trabajo: da soporte al \textit{pipeline} de captura y generación de datos aportando el fondo sobre
el que se entrenan y evalúan ambas políticas, pero ninguna conclusión del capítulo
\ref{Evaluación} depende de sus detalles. Se documenta aquí porque supuso un esfuerzo de
ingeniería considerable y porque uno de los problemas que planteó invalidó silenciosamente una
sesión completa de datos antes de detectarse. Se describe con el nivel de detalle necesario para
reproducirlo, no más; el propio proyecto de reconstrucción conserva la receta paso a paso.

\section{El activo y su integración en la escena}

La herramienta principal es \textbf{fVDB} \cite{williams2024fvdb}, el marco de NVIDIA para
aprendizaje profundo sobre datos espaciales dispersos, a través de su caja de herramientas de
captura de la realidad, \textit{fVDB Reality Capture} \cite{fvdb_reality_capture}, que implementa
la reconstrucción por \textit{Gaussian Splatting} \cite{kerbl2023gaussian} y la exportación a los
formatos que Isaac Sim lee. Las poses de cámara y la nube de puntos inicial de la que parte la
reconstrucción las produce \textbf{COLMAP} \cite{schonberger2016sfm}, la referencia en
\textit{Structure-from-Motion}, compilado con soporte CUDA para la arquitectura de las aceleradoras
de la estación. El resto es un puñado de guiones propios entre una etapa y la siguiente.

Hubo dos reconstrucciones de la oficina. La primera, el 3 de agosto de 2026, a partir de 101
fotografías de un teléfono móvil, sirvió para montar y validar el \textit{pipeline} entero,
incluida una malla de colisión que finalmente no hizo falta —la física de la escena es un plano de
suelo y la válvula, y el activo se usa solo como fondo visual—. La definitiva, al día siguiente,
partió de un \textbf{vídeo} de la misma oficina, y es la que llega al conjunto de datos. Sus
etapas y cifras:

\begin{enumerate}
	\item \textbf{Fotogramas.} Del vídeo se extraen 1089 fotogramas y se seleccionan 545 con una
	ventana deslizante que se queda con el más nítido de cada grupo —nitidez medida como la
	varianza del laplaciano—, en lugar de un umbral global, que descartaría bloques enteros justo
	donde la cámara giraba deprisa y rompería el emparejamiento secuencial.
	\item \textbf{Structure-from-Motion.} COLMAP con un único modelo de cámara compartido,
	extracción y emparejamiento de características en GPU, y reconstrucción incremental:
	\textbf{538 imágenes registradas} en un solo modelo, 142\,619 puntos tridimensionales, longitud
	media de traza 8,3 y error de reproyección medio de 1,18 píxeles; después, las imágenes se
	rectifican a un modelo de cámara sin distorsión.
	\item \textbf{Escala métrica y vertical.} COLMAP no conoce la escala. Un guion propio estima la
	vertical a partir de la orientación de las cámaras, ajusta un plano de suelo por RANSAC sobre
	los puntos más bajos —tras un recorte robusto por percentiles, sin el cual los puntos
	triangulados a través de ventanas y reflejos inflaban la escena hasta dimensiones absurdas—,
	rota la escena a $Z$ arriba con el suelo en $Z = 0$ y la escala fijando la altura media de las
	cámaras en 1,5\,m. Es una escala aproximada, con un error del orden del 5\,\%, suficiente para
	un fondo. Hacerlo \textbf{antes} de entrenar, y no corregirlo después sobre el activo
	exportado, es la decisión que evita toda una familia de problemas de matrices.
	\item \textbf{Reconstrucción.} \texttt{frgs reconstruct} sin normalización de escena —para
	conservar las unidades—, doscientas épocas sobre las 538 imágenes, 107\,600 pasos en unos
	treinta minutos en una GPU: \textbf{3\,962\,281 gaussianas}. Un guion propio las reduce después
	a 2\,144\,018 eliminando las de opacidad inferior a 0,005 —el refinado de la herramienta deja
	de podar a mitad del entrenamiento— y las de escala desproporcionada, que son pocas, suelen
	ser artefactos de zonas mal observadas y cuestan en render lo que un tile entero.
	\item \textbf{Exportación.} \texttt{frgs convert} al formato \textbf{NuRec} de Omniverse, en
	lugar del formato \textit{ParticleField} más reciente, porque sobre la misma reconstrucción
	NuRec se ve claramente mejor en Isaac Sim 6.0 y aquí solo importa cómo se ve desde la cámara
	del robot; la contrapartida es que NVIDIA lo va a retirar.
	\item \textbf{Envoltura e integración.} Una capa USD propia referencia el activo exportado y lo
	endereza y coloca —una traslación y una rotación de $-89$\textdegree{} en $X$, verificadas
	contra el dato: las 538 cámaras quedan a 1,50\,m de altura media y el suelo a 0,58\textdegree{}
	de la horizontal—, añade un plano de colisión de $10 \times 10$\,m en $Z = 0$ y una escena de
	física. El parche sobre la biblioteca de fondos del marco lo registra como fondo
	\texttt{office\_gs}, con una rotación de $+90$\textdegree{} en $Z$ para orientarlo respecto al
	robot y con variables de entorno para probar otro fichero, escala o pose sin tocar código.
\end{enumerate}

Cuatro trampas costaron más que el resto del proceso junto, y ninguna produjo un error.
\textbf{El exportador deja el activo tumbado}: escribe en la primitiva del volumen una «matriz de
conversión por defecto» pensada para datos con $Y$ arriba, que sobre un dato ya en $Z$ arriba
aplica una rotación que no hace falta; como esa matriz es su propia inversa, la corrección exacta
es ponerla a identidad, no buscar ángulos a ojo. \textbf{USD ignora los ajustes de render que el
activo trae consigo} cuando se referencia desde otra capa, de modo que el fichero de la tarea los
reaplica como ajustes del simulador cuando el fondo lleva la etiqueta \texttt{nurec}. \textbf{El
orden del cuaternión} de la pose en el marco es \textit{xyzw}, como dice el nombre del campo y al
contrario de lo que afirmaba una nota anterior; la identidad es $(0,0,0,1)$, y escribirla como
\textit{wxyz} tumba el fondo 89\textdegree. Se midió sobre el escenario ya compuesto con un guion
propio, tras tres intentos fallidos de predecirlo, porque hay tres capas de composición encima del
fichero de origen y con tres capas conviene medir. Y \textbf{el fondo heredado que existía al
empezar} —una oficina distinta, exportada con otra herramienta y de la que no se conservaban ni
las imágenes ni el espacio de trabajo— arrastraba una matriz transpuesta que USD lee como
proyectiva y que deformaba el activo en cuanto se le aplicaba una pose; es la razón de que se
capturara de nuevo en lugar de reutilizarlo.

\section{El procedimiento de reproducción sobre demostraciones grabadas}

El fondo no se renderiza durante la teleoperación: cuesta unos 39\,ms por paso de simulación, y
en las gafas eso es la diferencia entre poder agarrar un radio del volante y no poder. Las
demostraciones se graban en la escena diáfana y el fondo se aplica después, sin gafas y sin nadie
delante, con \texttt{rerender\_demos.py}: reproduce cada demostración con el fondo cargado, vuelve
a grabar la imagen de la cámara de la cabeza y escribe un nuevo HDF5 con la misma estructura. Cuesta
unos doce minutos por sesión de veinticinco demostraciones, y el fichero resultante ocupa
aproximadamente el doble que el original —110\,MB por demostración frente a 60— porque la oficina
comprime peor que un plano gris. Se comprobó que el fondo \textbf{no perturba la física}: la misma
demostración dura los mismos pasos con y sin él. El éxito se recalcula sobre la reproducción, no se
copia del original, y se verifica con \texttt{compare\_rerender.py} que cada demostración conserva
su número de pasos, cruza el umbral y ha cambiado de imagen.

\section{La divergencia de estado y su corrección}

La primera versión reproducía las acciones grabadas en lazo abierto, y \textbf{no reproducía la
demostración}. Se descubrió el 21 de agosto y había invalidado ya la salida de una sesión entera:
de veinticinco demostraciones que abrieron la válvula al teleoperar, solo dieciséis la abrían al
reproducirse, y la peor pasaba de 191,1\textdegree{} a 113,7\textdegree. Sobre la primera sesión,
frontal y sin agarre, la misma reproducción había perdido una de veinticinco y había parecido
correcta.

La opción \texttt{--validate\_states} de la herramienta, que compara en cada paso el estado
grabado con el reproducido, localizó la causa con exactitud: la divergencia aparece en el
\textbf{paso 1}, no gradualmente, y no en las posiciones articulares sino en las
\textbf{velocidades}, en los índices 14 y 18 —el cabeceo y el alabeo del tobillo derecho—. Con eso
la explicación es completa: el HDF5 almacena la acción de veintitrés dimensiones, que gobierna los
brazos y los canales de mando, pero \textbf{las piernas las gobierna la política de equilibrio en
lazo cerrado y no se reproducen: se vuelven a ejecutar}. Aterrizan en otro sitio desde el primer
paso, el tronco acaba ligeramente distinto, las muñecas agarran el radio por otro punto y al final
del episodio el volante se queda hasta 80\textdegree{} corto. Empujar de frente con la mano abierta
tolera esa deriva; agarrar un radio desde arriba, no.

La corrección no es hacer determinista la física, sino \textbf{no simular}: el HDF5 graba el
estado completo de la escena en cada paso, de modo que la herramienta lo impone tras cada
\texttt{step()} en lugar de dejar correr los controladores. Medido sobre las cincuenta
demostraciones de las sesiones 02 y 03: \textbf{50 de 50 reproducen, con una diferencia máxima de
1,1\textdegree{}} en el ángulo final, frente a los 81\textdegree{} anteriores, y sin fotogramas
negros. El residuo de 1,1\textdegree{} es un paso de física: el fotograma se renderiza dentro de
\texttt{step()}, justo antes de imponer el estado. Dos trampas de implementación: hay que usar
\texttt{env.scene.reset\_to()} y nunca \texttt{env.reset\_to()}, que además exporta un episodio y
\textbf{vuelve a sortear la disposición de la válvula en cada paso}; y en algunos caminos los
datos de la articulación llegan como \textit{arrays} de Warp, que no admiten indexado escalar y
fallan después de dos minutos de carga, por lo que se normalizan antes. La opción \texttt{--libre}
conserva el comportamiento antiguo, solo para comparar.

\section{Parámetros y configuración}

\begin{table*}[ht]
	\centering
	\caption{Parámetros del fondo fotorrealista y versiones de las herramientas}
	\label{tab:gs-parametros}
	\makebox[\textwidth]{
		\begin{tabular}{|l|p{9.6cm}|}
			\cline{1-2}
			Parámetro o herramienta & Valor \\ \hline
			\texttt{OFFICE\_GS\_USD} & Ruta del activo; por defecto el de la oficina de Madrid reconstruida a partir del vídeo \\
			\texttt{OFFICE\_GS\_ROT}, \texttt{\_POS}, \texttt{\_SCALE} & Pose y escala del fondo, cuaternión en orden \textit{xyzw}; por defecto $+90$\textdegree{} en $Z$, sin traslación, escala 1 \\
			\texttt{OFFICE\_GS\_LIGHT} & Intensidad de la luz de cúpula, 3000 por defecto. Solo se lee con el fondo fotorrealista; \textbf{todas} las reproducciones la dejan sin fijar, porque fijarla en una sesión y no en otra parte el conjunto de datos en dos exposiciones \\
			Ajustes de render & Los que trae el activo NuRec, reaplicados por la tarea como ajustes del simulador \\ \hline
			fVDB & \texttt{fvdb-core} 0.5.1 y \texttt{fvdb-reality-capture} 0.5.0, sobre PyTorch 2.11 con CUDA 12.8, en un entorno propio con Python 3.12 \\
			COLMAP & 4.2 de desarrollo, compilado con CUDA para la arquitectura \texttt{sm\_120}; el de la distribución no trae CUDA y sus opciones tienen otros nombres \\
			Isaac Sim & 6.0, con la transmisión de geometría activada, sin la cual las gaussianas se muestran como esferas grises \\ \hline
		\end{tabular}
	}
\end{table*}

Los guiones implicados, todos en el repositorio:

\begin{description}
	\item[\texttt{rerender\_demos.py}] la reproducción con el estado grabado impuesto.
	\item[\texttt{compare\_rerender.py}] la verificación de la reproducción contra el original.
	\item[\texttt{hdf5\_to\_video.py}] la revisión a ojo, sin arrancar el simulador.
	\item[\texttt{measure\_gs\_pose.py}] la medida de la pose real del fondo en el escenario compuesto.
	\item[\texttt{capture\_viewport.py}] la comprobación del render, antes de culpar a cualquier otra cosa.
\end{description}
